\makeatletter
% Parâmetros do documento (mesma interface do template dos relatórios)
\newcommand{\titulo}[1]{\gdef\@titulo{#1}}
\newcommand{\subtitulo}[1]{\gdef\@subtitulo{#1}}
\newcommand{\descricao}[1]{\gdef\@descricao{#1}}
\newcommand{\autores}[1]{\gdef\@autores{#1}}
\newcommand{\departamento}[1]{\gdef\@departamento{#1}}
\newcommand{\professor}[1]{\gdef\@professor{#1}}
\newcommand{\disciplina}[1]{\gdef\@disciplina{#1}}
\newcommand{\turma}[1]{\gdef\@turma{#1}}
\newcommand{\data}[1]{\gdef\@data{#1}}
% Campo a preencher antes da entrega (aparece em vermelho no PDF)
\newcommand{\pendente}[1]{\textcolor{red}{[#1]}}
% Figura que só é incluída se o arquivo existir em imgs/;
% caso contrário, desenha um quadro indicando o arquivo esperado.
% \figuraopcional{arquivo}{rótulo}{legenda}{o que a figura deve mostrar}
\newcommand{\figuraopcional}[4]{%
  \begin{figure}[H]\centering
  \IfFileExists{imgs/#1}{\includegraphics[width=0.9\linewidth]{imgs/#1}}%
  {\fbox{\parbox[c][4.5cm][c]{0.85\linewidth}{\centering\small
     \textcolor{red}{Inserir aqui: \texttt{imgs/#1}}\\[2mm]#4}}}%
  \caption{#3}\label{fig:#2}\end{figure}}
\newcommand{\figura}[3][0.85]{%
  \begin{figure}[H]\centering\includegraphics[width=#1\linewidth]{imgs/#2}
  \caption{#3}\label{fig:#2}\end{figure}}
\makeatother
